%% 2-2 マイナスの数と数直線（12ページ、p.26〜37）
%% ---- この章だけで使う絵の部品（頭文字 \bkiic…） ----
\makeatletter
% 数直線（この章ではいつも同じ形・同じ縮尺：-6〜6、1目もり 9 mm、0 が x=0）  \bkiicline{y}
\newcommand\bkiicline[1]{%
  \draw[bkInk, line width=0.9pt, -{Stealth[length=5pt]}] (-58.5,#1) -- (60.5,#1);
  \foreach \k in {-6,...,6} {\draw[bkInk, line width=0.6pt] (\k*9,#1-1.3) -- (\k*9,#1+1.3);
    \node[anchor=north, font=\footnotesize, text=bkInk, inner sep=0pt] at (\k*9,#1-2.6) {$\k$};}
  \draw[bkInk, line width=1.2pt] (0,#1-2) -- (0,#1+2);}
% 数直線の上の点  \bkiicdot{x}{y}{色}
\newcommand\bkiicdot[3]{\fill[#3] (#1*9,#2) circle (1.5);}
% 数直線の上を進む矢印  \bkiicmove{はじめ}{おわり}{y}{色}{ラベル}
\newcommand\bkiicmove[5]{\draw[bk arrow, #4, line width=1.4pt] (#1*9,#3) -- node[above=0.8mm, font=\sffamily\small, text=#4!80!black, inner sep=0pt] {#5} (#2*9,#3);}
% 温度計（目もりの数字はマイナス記号で書く）  \bkiicthermo{(x,y)}{1度の長さ}{いちばん下}{いちばん上}{いまの温度}
%   (x,y) は 0 ℃ の目もりの位置。数字は5度ごとに左へ
\newcommand\bkiicthermo[5]{\begin{scope}[shift={#1}]
  \fill[white, draw=bkGray, line width=0.8pt, rounded corners=2.5mm] (-4,#3*#2-9) rectangle (4,#4*#2+4);
  \fill[bkRed] (0,#3*#2-5.5) circle (2.6);
  \fill[bkRed] (-0.9,#3*#2-4) rectangle (0.9,#5*#2);
  \foreach \t in {#3,...,#4} {\pgfmathtruncatemacro\bkiic@lab{mod(\t,5)==0}
    \ifnum\bkiic@lab=1 \draw[bkInk, line width=0.7pt] (1.4,\t*#2) -- (4,\t*#2);
      \node[anchor=east, font=\footnotesize, text=bkInk, inner sep=0pt] at (-4.8,\t*#2) {$\t$};
    \else\draw[bkInk!80, line width=0.4pt] (2.4,\t*#2) -- (4,\t*#2);\fi}
\end{scope}}
\makeatother

\bkchapter{2}{2}{マイナスの\R{数}{かず}と\R{数直線}{すう|ちょく|せん}}{\bklevelrange{1}{2}}{マイナスの\R{数}{かず}と\R{数直線}{すう|ちょく|せん}}
\bkmichishirube{1-3,1-6,1-11}{2-3,3-2,4-10,5-5}

\begin{goal}
\begin{itemize}[label={\color{bkTheme}●}, leftmargin=1.6\zw, itemsep=0pt, topsep=1pt]
\item \bklevel{1}\enspace 0より\R{小}{ちい}さい\R{数}{かず}を\R{負}{ふ}の\R{数}{かず}で\R{表}{あらわ}し、\R{数直線}{すう|ちょく|せん}にかける。
\item \bklevel{1}\enspace \R{不等号}{ふ|とう|ごう}（$<$・$>$・$\leqq$・$\geqq$）で、\R{数}{かず}の\R{大}{おお}きさをくらべられる。
\item \bklevel{2}\enspace \R{正}{せい}の\R{数}{かず}・\R{負}{ふ}の\R{数}{かず}の ＋ － × ÷ の\R{計算}{けい|さん}ができる。
\end{itemize}
\end{goal}

\begin{hatena}
\noindent\begin{minipage}[c]{0.66\linewidth}
\R{冬}{ふゆ}の\R{朝}{あさ}、\R{校庭}{こう|てい}の\R{温度計}{おん|ど|けい}を\R{見}{み}ると、\R{赤}{あか}い\R{液}{えき}の\R{先}{さき}が、0の\R{目}{め}もりより3\R{目}{め}もり\R{下}{した}にありました。\R{昼}{ひる}には、0より3\R{目}{め}もり\R{上}{うえ}になりました。

0\,℃ は、\R{温度}{おん|ど}が「ない」という\R{意味}{い|み}かな？\quad 0より\R{下}{した}の\R{温度}{おん|ど}は、どう\R{書}{か}けばいいのかな？
\end{minipage}\hfill
\begin{minipage}[c]{0.31\linewidth}\centering
\begin{tikzpicture}[x=1mm, y=1mm]
  \bkiicthermo{(10,0)}{2.4}{-10}{10}{-3}
  \bkiicthermo{(30,0)}{2.4}{-10}{10}{3}
  \node[anchor=north, font=\sffamily\footnotesize, inner sep=0pt] at (10,-34.5) {\R{朝}{あさ}};
  \node[anchor=north, font=\sffamily\footnotesize, inner sep=0pt] at (30,-34.5) {\R{昼}{ひる}};
\end{tikzpicture}
\end{minipage}
\end{hatena}

\R{第}{だい}1\R{巻}{かん}の\R{最後}{さい|ご}のカード「0より\R{小}{ちい}さいかずってある？」の\R{答}{こた}えも、この\R{章}{しょう}でわかります。

\clearpage
%% ---------------- p.27 やってみよう・思い出そう ----------------
\begin{yattemiyou}[\R{氷}{こおり}に\R{塩}{しお}をかけて、\R{温度}{おん|ど}をはかろう]
\noindent\begin{minipage}[c]{0.66\linewidth}
\begin{enumerate}[label={\bkcircled[bkGreen]{\arabic*}}, leftmargin=2.4\zw, itemsep=1pt, topsep=0pt]
\item コップに\R{氷}{こおり}と\R{少}{すこ}しの\R{水}{みず}を\R{入}{い}れて、\R{温度}{おん|ど}をはかろう。
\item \R{氷}{こおり}に\R{塩}{しお}をたくさんかけて、よくまぜよう。
\item 1\R{分}{ぷん}ごとに\R{温度}{おん|ど}を\R{読}{よ}み、ノートの\R{表}{ひょう}に\R{書}{か}こう。0の\R{目}{め}もりより\R{下}{した}になったら、どう\R{書}{か}けばいいかな？
\end{enumerate}
{\footnotesize\color{bkGray}\R{冷}{つめ}たくなるので、コップを\R{長}{なが}くさわらないようにしよう。}
\end{minipage}\hfill
\begin{minipage}[c]{0.31\linewidth}\centering
\begin{tikzpicture}[x=1mm, y=1mm]
  % コップと氷と温度計
  \fill[bkBlue!12] (0.8,0.8) rectangle (21.2,20);
  \foreach \x/\y in {4/4, 10/3.5, 16/4.5, 6.5/10, 13/9.5, 18/11, 4/15.5, 11/16} {\filldraw[fill=white, draw=bkBlue!50, line width=0.5pt, rotate around={15:(\x,\y)}] (\x-2.4,\y-2.2) rectangle (\x+2.4,\y+2.2);}
  \draw[bkInk!80, line width=1pt, line join=round] (0,26) -- (0,0) -- (22,0) -- (22,26);
  \fill[white, draw=bkGray, line width=0.7pt, rounded corners=1.5mm] (12,3) rectangle (16.5,38);
  \fill[bkRed] (14.25,6) circle (1.6);
  \fill[bkRed] (13.7,6) rectangle (14.8,14);
  \foreach \yy in {16,20,...,36} \draw[bkInk, line width=0.4pt] (15,\yy) -- (16.5,\yy);
\end{tikzpicture}
\end{minipage}

\end{yattemiyou}

\begin{omoidasou}
\bkrecall{1-6}{\R{数直線}{すう|ちょく|せん}は、\R{右}{みぎ}へ\R{行}{い}くほど\R{大}{おお}きい。たし\R{算}{ざん}は\R{右}{みぎ}へ、ひき\R{算}{ざん}は\R{左}{ひだり}へ\R{進}{すす}む。}
\bkrecall{1-5}{「かわった\R{量}{りょう}」は、あと－まえ で\R{求}{もと}める。}
\bkrecall{1-12}{かけ\R{算}{ざん}は「1つ\R{分}{ぶん}の\R{数}{かず}×いくつ\R{分}{ぶん}」。わり\R{算}{ざん}はかけ\R{算}{ざん}の\R{逆}{ぎゃく}（1-14）。}
\end{omoidasou}

%% ---------------- 数学の道具 ----------------
\bksection{doguu}{\R{数学}{すう|がく}の\R{道具}{どう|ぐ}}

\bksubsection{1}{0より\R{小}{ちい}さい\R{数}{かず}}

\leavevmode\label{bk:2-2:fu}\bkyattemiyouref では、\R{塩}{しお}をかけた\R{氷}{こおり}が、0の\R{目}{め}もりより10\R{目}{め}もりほど\R{下}{した}まで\R{冷}{ひ}えたはずです。

0\,℃ は、\R{温度}{おん|ど}が「ない」という\R{意味}{い|み}ではありません。\R{水}{みず}がこおる\R{温度}{おん|ど}は、だれでもたしかめられます。そこで、この\R{温度}{おん|ど}を、\R{目}{め}もりの\R{数}{かず}を\R{決}{き}めるもとの\R{場所}{ば|しょ}（\textbf{\R{基準}{き|じゅん}}）にして、0と\R{決}{き}めたのです。0\,℃ より\R{寒}{さむ}い\R{日}{ひ}も、ちゃんと\R{温度}{おん|ど}はあります。

\R{基準}{き|じゅん}より\R{低}{ひく}い\R{温度}{おん|ど}は、マイナス（$-$）をつけて \qty{-10}{\degreeCelsius} と\R{書}{か}き、「マイナス10\R{度}{ど}」と\R{読}{よ}みます（\R{天気予報}{てん|き|よ|ほう}では「\R{氷点下}{ひょう|てん|か}10\R{度}{ど}」ともいいます）。\bkhatenaref の\R{朝}{あさ}の\R{温度}{おん|ど}は \qty{-3}{\degreeCelsius}、\R{昼}{ひる}は \qty{3}{\degreeCelsius} です。どちらも0から3\R{目}{め}もりはなれていますが、\R{向}{む}きが\R{反対}{はん|たい}です。

\noindent\begin{minipage}[c]{0.50\linewidth}\centering
\begin{tikzpicture}[x=1mm, y=1mm]
  % 建物の断面（1 m を 2.5 mm で。1階の高さ 4 m）
  \fill[bkYellowBg] (-6,-24) rectangle (38,0);
  \fill[white, draw=bkInk!80, line width=0.8pt] (0,-20) rectangle (30,20);
  \foreach \yy in {-10,0,10} {\draw[bkInk!80, line width=1.2pt] (0,\yy) -- (30,\yy);}
  \draw[bkStem, line width=1.4pt] (-6,0) -- (0,0) (30,0) -- (38,0);
  \node[font=\sffamily\footnotesize, inner sep=0pt] at (15,15) {2\R{階}{かい}};
  \node[font=\sffamily\footnotesize, inner sep=0pt] at (15,5) {1\R{階}{かい}};
  \node[font=\sffamily\footnotesize, inner sep=0pt] at (15,-5) {\R{地下}{ち|か}1\R{階}{かい}};
  \node[font=\sffamily\footnotesize, inner sep=0pt] at (15,-15) {\R{地下}{ち|か}2\R{階}{かい}};
  \node[anchor=south east, font=\sffamily\scriptsize, text=bkStem, inner sep=0pt] at (-1,1) {\R{地面}{じ|めん}};
  % 高さの数直線
  \draw[bkInk, line width=0.9pt, -{Stealth[length=4pt]}] (50,-24) -- (50,25);
  \node[anchor=south, font=\sffamily\footnotesize, inner sep=0pt] at (50,26.5) {\R{高}{たか}さ（m）};
  \foreach \yy/\l in {20/8, 10/4, 0/0, -10/-4, -20/-8} {
    \draw[bkInk, line width=0.6pt] (48.8,\yy) -- (51.2,\yy);
    \node[anchor=west, font=\footnotesize, inner sep=0pt] at (53,\yy) {$\l$};}
  \foreach \yy in {10,0,-10,-20} {\draw[bkGray, line width=0.4pt, densely dashed] (31.5,\yy) -- (47.5,\yy);}
\end{tikzpicture}
\end{minipage}\hfill
\begin{minipage}[c]{0.47\linewidth}
\R{建物}{たて|もの}でも\R{同}{おな}じです。1\R{階}{かい}の\R{床}{ゆか}（\R{地面}{じ|めん}）を\R{基準}{き|じゅん}の0にします。1つの\R{階}{かい}の\R{高}{たか}さを \qty{4}{m} とすると、2\R{階}{かい}の\R{床}{ゆか}は \qty{4}{m}、\R{地下}{ち|か}1\R{階}{かい}の\R{床}{ゆか}は \qty{-4}{m}、\R{地下}{ち|か}2\R{階}{かい}の\R{床}{ゆか}は \qty{-8}{m}。エレベーターのボタンの「B1」は、\R{地下}{ち|か}1\R{階}{かい}のことです。
\end{minipage}

\par\vspace{2mm}
\R{温度}{おん|ど}や\R{建物}{たて|もの}のほかにも、\R{基準}{き|じゅん}を0にして、\R{上}{うえ}を＋、\R{下}{した}を－で\R{表}{あらわ}す\R{量}{りょう}があります。

\begin{bkcenter}
\begin{bktable}{colspec={Q[l,m,wd=28mm] Q[l,m,wd=34mm] Q[c,m,wd=26mm] Q[c,m,wd=26mm]}, rowsep=1pt, cells={font=\small}, row{1}={font=\sffamily\bfseries\small}}
 & \R{基準}{き|じゅん}（0） & \R{基準}{き|じゅん}より\R{上}{うえ} & \R{基準}{き|じゅん}より\R{下}{した} \\
\R{気温}{き|おん} & \R{水}{みず}がこおる\R{温度}{おん|ど} & \qty{3}{\degreeCelsius} & \qty{-3}{\degreeCelsius} \\
\R{土地}{と|ち}の\R{高}{たか}さ & \R{海面}{かい|めん} & \R{富士山}{ふ|じ|さん} \qty{3776}{m} & \R{死海}{し|かい} \R{約}{やく}\qty{-430}{m} \\
\R{建物}{たて|もの}の\R{高}{たか}さ & \R{地面}{じ|めん} & 2\R{階}{かい}の\R{床}{ゆか} \qty{4}{m} & \R{地下}{ち|か}1\R{階}{かい}の\R{床}{ゆか} \qty{-4}{m} \\
\R{貯金}{ちょ|きん}のふえ\R{方}{かた} & \R{今}{いま}の\R{金額}{きん|がく} & \R{入}{はい}る：$+500$\R{円}{えん} & \R{使}{つか}う：$-300$\R{円}{えん} \\
\end{bktable}
\end{bkcenter}

0より\R{大}{おお}きい\R{数}{かず}を\textbf{\R{正}{せい}の\R{数}{かず}}、0より\R{小}{ちい}さい\R{数}{かず}を\textbf{\R{負}{ふ}の\R{数}{かず}}といいます。\R{正}{せい}の\R{数}{かず}は $+3$（プラス3）のように＋をつけることもありますが、ふつうは\R{省}{はぶ}いて3と\R{書}{か}きます。0は、\R{正}{せい}の\R{数}{かず}でも\R{負}{ふ}の\R{数}{かず}でもありません。\bkidx{せいのかず}{\R{正}{せい}の\R{数}{かず}}\bkidx{ふのかず}{\R{負}{ふ}の\R{数}{かず}}

\bksubsection{2}{\R{数直線}{すう|ちょく|せん}を\R{左}{ひだり}へのばす}

\leavevmode\R{温度計}{おん|ど|けい}を\R{右}{みぎ}へたおすと、\R{上}{うえ}が\R{右}{みぎ}、\R{下}{した}が\R{左}{ひだり}になり、\R{数直線}{すう|ちょく|せん}になります。\R{負}{ふ}の\R{数}{かず}は、0の\R{左}{ひだり}にならびます。1\R{目}{め}もりの\R{長}{なが}さは、0の\R{右}{みぎ}も\R{左}{ひだり}も\R{同}{おな}じです。

\begin{bkfigure}
\begin{tikzpicture}[x=1mm, y=1mm]
  \fill[bkBlue!12] (-56,-1) rectangle (-1,7);
  \fill[bkOrange!14] (1,-1) rectangle (56,7);
  \node[font=\sffamily\small, text=bkBlue!80!black] at (-38,4.2) {\R{負}{ふ}の\R{数}{かず}};
  \node[font=\sffamily\small, text=bkOrange!80!black] at (28,4.2) {\R{正}{せい}の\R{数}{かず}};
  \bkiicline{0}
  \bkiicdot{-5}{0}{bkBlue}\bkiicdot{-2}{0}{bkBlue}\bkiicdot{-2.5}{0}{bkPurple}
  \node[anchor=south, font=\footnotesize, text=bkPurple, inner sep=0pt] at (-22.5,2.2) {$-2.5$};
  \draw[bk arrow, bkGray, line width=1pt] (-36,-10.5) -- node[below, font=\sffamily\small, text=bkInk, inner sep=1.5mm] {\R{小}{ちい}さい} (-56,-10.5);
  \draw[bk arrow, bkGray, line width=1pt] (36,-10.5) -- node[below, font=\sffamily\small, text=bkInk, inner sep=1.5mm] {\R{大}{おお}きい} (56,-10.5);
\end{tikzpicture}
\bkfigcaption{\R{右}{みぎ}へ\R{行}{い}くほど\R{大}{おお}きく、\R{左}{ひだり}へ\R{行}{い}くほど\R{小}{ちい}さい。}
\end{bkfigure}

$-5$ と $-2$ では、どちらが\R{大}{おお}きいかな？\quad $-5$ は $-2$ より\R{左}{ひだり}にあるので、$-5$ のほうが\R{小}{ちい}さい\R{数}{かず}です。\R{温度}{おん|ど}で\R{考}{かんが}えても、\qty{-5}{\degreeCelsius} のほうが \qty{-2}{\degreeCelsius} より\R{寒}{さむ}い（\R{温度}{おん|ど}が\R{低}{ひく}い）ですね。\R{負}{ふ}の\R{数}{かず}では、0から\R{遠}{とお}いほど\R{小}{ちい}さくなります。

\R{小数}{しょう|すう}の\R{負}{ふ}の\R{数}{かず}もあります。$-2.5$ は、$-2$ と $-3$ のちょうどまん\R{中}{なか}です（\R{図}{ず}のむらさきの\R{点}{てん}）。

\bksubsection{3}{\R{大}{おお}きさをくらべる\R{記号}{き|ごう}：\R{不等号}{ふ|とう|ごう}}

\leavevmode\label{bk:2-2:futougou}「$-5$ は $-2$ より\R{小}{ちい}さい」を、\R{記号}{き|ごう}で $-5<-2$ と\R{書}{か}きます。「$3$ は $-1$ より\R{大}{おお}きい」は $3>-1$ です。\R{記号}{き|ごう}の\R{口}{くち}が\R{開}{ひら}いているほうに、\R{大}{おお}きい\R{数}{かず}を\R{書}{か}きます。$<$、$>$ を\textbf{\R{不等号}{ふ|とう|ごう}}といいます。\bkidx{ふとうごう}{\R{不等号}{ふ|とう|ごう}}$-5<-2$ は「$-5$ は $-2$ より\R{小}{ちい}さい」と\R{読}{よ}みます。

\noindent\begin{minipage}[c]{0.40\linewidth}\centering
\begin{tikzpicture}[x=1mm, y=1mm]
  \fill[bkGray!70] (21.5,-4) rectangle (24.5,0);
  \fill[bkYellowBg, draw=bkYellow!80!black, line width=0.8pt, rounded corners=2mm] (0,0) rectangle (46,17);
  \node[font=\sffamily\bfseries\small, text=bkInk, align=center] at (23,8.5) {\R{身長}{しん|ちょう} \qty{120}{cm} \R{以上}{い|じょう}の\\\R{人}{ひと}が\R{乗}{の}れます};
\end{tikzpicture}
\end{minipage}\hfill
\begin{minipage}[c]{0.57\linewidth}
\qty{120}{cm} ちょうどの\R{人}{ひと}は\R{乗}{の}れるかな？\quad「120\R{以上}{い|じょう}」は、120と、120より\R{大}{おお}きい\R{数}{かず}です。120ちょうどもふくむので、\R{乗}{の}れます。
\end{minipage}

\par\vspace{2mm}
「ちょうど」をふくむかどうかで、ことばと\R{記号}{き|ごう}がかわります。$\geqq$ は $>$ と、「\R{等}{ひと}しい」の $=$ を\R{合}{あ}わせた\R{形}{かたち}です。\R{下}{した}の\R{線}{せん}が「ちょうどその\R{数}{かず}もふくむ」という\R{印}{しるし}です。

\begin{bkcenter}
\begin{bktable}{colspec={Q[l,m,wd=58mm] Q[c,m,wd=20mm] Q[l,m,wd=30mm]}, rowsep=0.5pt, cells={font=\small}, row{1}={font=\sffamily\bfseries\small}}
ことば & \R{記号}{き|ごう} & 3をふくむ？ \\
3\R{以上}{い|じょう}（3と、3より\R{大}{おお}きい） & $\text{□}\geqq3$ & ふくむ \\
3\R{以下}{い|か}（3と、3より\R{小}{ちい}さい） & $\text{□}\leqq3$ & ふくむ \\
3より\R{大}{おお}きい（3をこえる） & $\text{□}>3$ & ふくまない \\
3\R{未満}{み|まん}（3より\R{小}{ちい}さい） & $\text{□}<3$ & ふくまない \\
\end{bktable}
\end{bkcenter}

\R{数直線}{すう|ちょく|せん}にかくと、\R{下}{した}のようになります。●はその\R{数}{かず}をふくむ、○はふくまない、という\R{印}{しるし}です。

\begin{bkfigure}
\begin{tikzpicture}[x=1mm, y=1mm]
  \foreach \yy/\lab in {0/{3\R{以上}{い|じょう}\enspace $\text{□}\geqq3$}, -14/{3\R{未満}{み|まん}\enspace $\text{□}<3$}, -28/{0より\R{大}{おお}きく6より\R{小}{ちい}さい\\$0<\text{□}<6$}} {
    \draw[bkInk, line width=0.8pt, -{Stealth[length=4pt]}] (-12,\yy) -- (84,\yy);
    \foreach \k in {-1,...,8} {\draw[bkInk, line width=0.5pt] (\k*10,\yy-1.2) -- (\k*10,\yy+1.2);
      \node[anchor=north, font=\footnotesize, inner sep=0pt] at (\k*10,\yy-2.4) {$\k$};}
    \node[anchor=east, font=\sffamily\footnotesize, align=right, inner sep=0pt] at (-15,\yy) {\lab};}
  % 3以上
  \draw[bkOrange, line width=2.2pt] (30,0) -- (81,0);
  \fill[bkOrange] (30,0) circle (1.8);
  % 3未満
  \draw[bkBlue, line width=2.2pt] (-11,-14) -- (28.2,-14);
  \draw[bkBlue, line width=1pt, fill=white] (30,-14) circle (1.8);
  % 0<□<6
  \draw[bkGreen!80!black, line width=2.2pt] (1.8,-28) -- (58.2,-28);
  \draw[bkGreen!80!black, line width=1pt, fill=white] (0,-28) circle (1.8) (60,-28) circle (1.8);
\end{tikzpicture}
\bkfigcaption{●はふくむ、○はふくまない。}
\end{bkfigure}

…、$-2$、$-1$、0、1、2、…のような\R{数}{かず}を\textbf{\R{整数}{せい|すう}}といいます。\bkidx{せいすう}{\R{整数}{せい|すう}}0.5 や $-2.5$ や $\dfrac{1}{2}$ のように、\R{整数}{せい|すう}と\R{整数}{せい|すう}のあいだにある\R{数}{かず}は、\R{整数}{せい|すう}ではありません。

$0<\text{□}<6$ は、「$0<\text{□}$（□は0より\R{大}{おお}きい）」と「$\text{□}<6$（□は6より\R{小}{ちい}さい）」の2つを、まん\R{中}{なか}の□でつないで、1\R{行}{ぎょう}にしたものです。□の\R{数}{かず}は、0と6のあいだにあります（\R{上}{うえ}の\R{図}{ず}のいちばん\R{下}{した}）。

\R{物理}{ぶつ|り}の\R{本}{ほん}では、□のかわりに $x$ や $a$ のような\R{文字}{も|じ}を\R{使}{つか}います（\R{文字}{も|じ}は2-11でくわしく\R{学}{まな}びます）。$0<x<6$ は「$x$ は0より\R{大}{おお}きく6より\R{小}{ちい}さい」、$a>0$ は「$a$ は\R{正}{せい}」、$a<0$ は「$a$ は\R{負}{ふ}」と\R{読}{よ}みます。\bkidx{いじょういかみまん}{\R{以上}{い|じょう}・\R{以下}{い|か}・\R{未満}{み|まん}}

\bksubsection{4}{＋の\R{向}{む}きは、\R{自分}{じ|ぶん}で\R{決}{き}める}

\leavevmode\label{bk:2-2:muki}\R{東西}{とう|ざい}にのびる\R{道}{みち}で、ソラは\R{東}{ひがし}へ \qty{3}{m}、ミナは\R{西}{にし}へ \qty{3}{m} \R{歩}{ある}きました。\R{東}{ひがし}を＋と\R{決}{き}めると、ソラは $+\qty{3}{m}$、ミナは $\qty{-3}{m}$ です。\R{西}{にし}は\R{東}{ひがし}の\R{反対}{はん|たい}の\R{向}{む}きなので、マイナスをつけるのです。

\begin{bkfigure}
\begin{tikzpicture}[x=1mm, y=1mm]
  \node[anchor=south west, font=\sffamily\small, inner sep=0pt] at (-58.5,2.5) {\R{西}{にし}};
  \node[anchor=south east, font=\sffamily\small, inner sep=0pt] at (60.5,2.5) {\R{東}{ひがし}（＋）};
  \bkiicline{0}
  \bkiicmove{0}{3}{6}{bkOrange}{ソラ $+3$}
  \bkiicmove{0}{-3}{6}{bkBlue}{ミナ $-3$}
\end{tikzpicture}
\bkfigcaption{\R{東}{ひがし}を＋と\R{決}{き}めると、\R{西}{にし}\R{向}{む}きはマイナスで\R{表}{あらわ}す。}
\end{bkfigure}

\R{西}{にし}を＋と\R{決}{き}めれば、ソラが $-3$、ミナが $+3$ になります。どちらを＋にしてもかまいません。でも、とちゅうで\R{変}{か}えると、\R{答}{こた}えがめちゃくちゃになります（\pageref{bk:2-2:tochu}ページのまちがえやすいところ）。

マイナスは「\R{反対}{はん|たい}の\R{向}{む}き」を\R{表}{あらわ}す。＋の\R{向}{む}きは、はじめに\R{決}{き}めて\R{変}{か}えない。これが\R{大事}{だい|じ}なきまりです。

\bksubsection{5}{たし\R{算}{ざん}は、\R{数直線}{すう|ちょく|せん}の\R{上}{うえ}の\R{動}{うご}き}

\leavevmode\label{bk:2-2:keisan}\R{朝}{あさ}の\R{気温}{き|おん}は \qty{-3}{\degreeCelsius}。\R{昼}{ひる}までに \qty{5}{\degreeCelsius} \R{上}{あ}がり、\R{夜}{よる}までに \qty{4}{\degreeCelsius} \R{下}{さ}がりました。\R{昼}{ひる}と\R{夜}{よる}の\R{気温}{き|おん}は\R{何}{なん}℃ かな？\quad \R{数直線}{すう|ちょく|せん}の\R{上}{うえ}で\R{動}{うご}いてみよう。

\begin{bkfigure}
\begin{tikzpicture}[x=1mm, y=1mm]
  \bkiicline{0}
  \bkiicdot{-3}{0}{bkInk}
  \bkiicmove{-3}{2}{5}{bkOrange}{$+5$（\R{右}{みぎ}へ5）}
  \bkiicmove{2}{-2}{13}{bkBlue}{$+(-4)$（\R{左}{ひだり}へ4）}
\end{tikzpicture}
\bkfigcaption{\R{正}{せい}の\R{数}{かず}をたすと\R{右}{みぎ}へ、\R{負}{ふ}の\R{数}{かず}をたすと\R{左}{ひだり}へ\R{進}{すす}む。}
\end{bkfigure}

\noindent{\sffamily\bfseries\color{bkDeep}①\enspace\R{正}{せい}の\R{数}{かず}をたす}\par
\qty{5}{\degreeCelsius} \R{上}{あ}がることは、$+5$ です。たし\R{算}{ざん}は\R{右}{みぎ}へ\R{進}{すす}むことでした（1-6）。$(-3)+5$ は、$-3$ から\R{右}{みぎ}へ5\R{進}{すす}んで2。\R{昼}{ひる}の\R{気温}{き|おん}は \qty{2}{\degreeCelsius} です。

\par\vspace{1mm}
\noindent{\sffamily\bfseries\color{bkDeep}②\enspace\R{負}{ふ}の\R{数}{かず}をたす}\par
\qty{4}{\degreeCelsius} \R{下}{さ}がることは、「\R{上}{あ}がる」の\R{反対}{はん|たい}なので、$-4$ と\R{表}{あらわ}します。それをたすので、$2+(-4)$ と\R{書}{か}きます。「$-4$ をたす」は「\R{右}{みぎ}へ $-4$ \R{進}{すす}む」こと。マイナスは\R{反対}{はん|たい}の\R{向}{む}きなので、「\R{左}{ひだり}へ4\R{進}{すす}む」ことです。2から\R{左}{ひだり}へ4\R{進}{すす}んで $-2$。\R{夜}{よる}の\R{気温}{き|おん}は \qty{-2}{\degreeCelsius} です。

\R{負}{ふ}の\R{数}{かず}に\R{負}{ふ}の\R{数}{かず}をたすときも\R{同}{おな}じです。$(-3)+(-2)$ は、$-3$ から\R{左}{ひだり}へ2\R{進}{すす}んで $-5$ です。

\par\vspace{1mm}
\noindent{\sffamily\bfseries\color{bkDeep}③\enspace お\R{金}{かね}でたしかめる}\par
おこづかいを500\R{円}{えん}もらい（$+500$）、300\R{円}{えん}\R{使}{つか}いました（$-300$）。お\R{金}{かね}は、$500+(-300)=200$ で、200\R{円}{えん}ふえています。500\R{円}{えん}から300\R{円}{えん}\R{使}{つか}うと200\R{円}{えん}\R{残}{のこ}るので、\R{合}{あ}っていますね。\R{負}{ふ}の\R{数}{かず}をたすと、\R{数}{かず}は\R{小}{ちい}さくなります。

\begin{bkcenter}
\begin{bktable}{colspec={Q[l,m,wd=40mm] Q[l,m,wd=36mm] Q[l,m,wd=36mm]}, rowsep=1pt, cells={font=\small}, row{1}={font=\sffamily\bfseries\small}}
 & \R{数直線}{すう|ちょく|せん}の\R{上}{うえ}で & \R{例}{れい} \\
\R{正}{せい}の\R{数}{かず}をたす & \R{右}{みぎ}へ\R{進}{すす}む & $(-3)+5=2$ \\
\R{負}{ふ}の\R{数}{かず}をたす & \R{左}{ひだり}へ\R{進}{すす}む & $2+(-4)=-2$ \\
\end{bktable}
\bktabcaption{たし\R{算}{ざん}のまとめ}
\end{bkcenter}

\bksubsection{6}{ひき\R{算}{ざん}は、\R{反対}{はん|たい}の\R{数}{かず}をたすこと}

\leavevmode ひき\R{算}{ざん}は\R{左}{ひだり}へ\R{進}{すす}むことでした（1-6）。$2-6$ は、2から\R{左}{ひだり}へ6\R{進}{すす}んで $-4$ です。0より\R{左}{ひだり}にも\R{数}{かず}があるので、\R{小}{ちい}さい\R{数}{かず}から\R{大}{おお}きい\R{数}{かず}をひくこともできます。

では、$2-(-3)$ のように、\R{負}{ふ}の\R{数}{かず}をひくときはどうなるでしょう。「$-3$ をひく」は「\R{左}{ひだり}へ $-3$ \R{進}{すす}む」こと。マイナスは\R{反対}{はん|たい}の\R{向}{む}きなので、「\R{右}{みぎ}へ3\R{進}{すす}む」ことです。2から\R{右}{みぎ}へ3\R{進}{すす}んで5です。

\begin{bkfigure}
\begin{tikzpicture}[x=1mm, y=1mm]
  \bkiicline{0}
  \bkiicdot{2}{0}{bkInk}
  \bkiicmove{2}{-4}{5}{bkBlue}{$2-6$：\R{左}{ひだり}へ6}
  \bkiicmove{2}{5}{14.5}{bkOrange}{$2-(-3)$：\R{右}{みぎ}へ3}
\end{tikzpicture}
\bkfigcaption{\R{正}{せい}の\R{数}{かず}をひくと\R{左}{ひだり}へ、\R{負}{ふ}の\R{数}{かず}をひくと\R{右}{みぎ}へ\R{進}{すす}む。}
\end{bkfigure}

\R{本当}{ほん|とう}に $2-(-3)=5$ になるのか、\R{別}{べつ}の\R{方法}{ほう|ほう}でもたしかめましょう。2から、ひく\R{数}{かず}を1ずつ\R{小}{ちい}さくしていきます。

\begin{bkcenter}
\begin{bktable}{colspec={Q[l,m,wd=9mm] *{7}{Q[c,m,wd=15.5mm]}}, rowsep=1pt, colsep=1.5pt, cells={font=\small}, column{1}={font=\sffamily\bfseries\small}}
\R{式}{しき} & $2-3$ & $2-2$ & $2-1$ & $2-0$ & $2-(-1)$ & $2-(-2)$ & $2-(-3)$ \\
\R{答}{こた}え & $-1$ & 0 & 1 & 2 & {\color{bkRed}\bfseries 3} & {\color{bkRed}\bfseries 4} & {\color{bkRed}\bfseries 5} \\
\end{bktable}
\bktabcaption{ひく\R{数}{かず}が1\R{小}{ちい}さくなると、\R{答}{こた}えは1\R{大}{おお}きくなる。}
\end{bkcenter}

\R{左}{ひだり}の4つは、\R{今}{いま}までの\R{計算}{けい|さん}でわかります。ひく\R{数}{かず}が1\R{小}{ちい}さくなるたびに、ひかれる\R{量}{りょう}がへるので、\R{答}{こた}えは1ずつ\R{大}{おお}きくなっています。このならびを\R{右}{みぎ}へつづけると、\R{赤}{あか}い\R{数}{かず}になります。$2-(-3)=5$ で、\R{数直線}{すう|ちょく|せん}と\R{同}{おな}じ\R{答}{こた}えです。

\noindent\begin{minipage}[c]{0.50\linewidth}\centering
\begin{tikzpicture}[x=1mm, y=1mm]
  \draw[bkInk, line width=0.9pt, -{Stealth[length=4pt]}] (-33,0) -- (34,0);
  \foreach \k in {-3,...,3} {\draw[bkInk, line width=0.6pt] (\k*9,-1.3) -- (\k*9,1.3);
    \node[anchor=north, font=\footnotesize, inner sep=0pt] at (\k*9,-2.6) {$\k$};}
  \draw[bkInk, line width=1.2pt] (0,-2) -- (0,2);
  \fill[bkBlue] (-27,0) circle (1.5);\fill[bkOrange] (27,0) circle (1.5);
  \draw[bkBlue, line width=0.6pt, {Stealth[length=3pt]}-{Stealth[length=3pt]}] (-1.5,5) -- (-25.5,5);
  \draw[bkOrange, line width=0.6pt, {Stealth[length=3pt]}-{Stealth[length=3pt]}] (1.5,5) -- (25.5,5);
  \node[anchor=south, font=\sffamily\footnotesize, inner sep=0pt] at (-13.5,6.5) {3};
  \node[anchor=south, font=\sffamily\footnotesize, inner sep=0pt] at (13.5,6.5) {3};
\end{tikzpicture}
\end{minipage}\hfill
\begin{minipage}[c]{0.47\linewidth}
3と $-3$ のように、0をはさんで\R{反対}{はん|たい}がわの、0から\R{同}{おな}じ\R{距離}{きょ|り}にある\R{数}{かず}を、\textbf{\R{反対}{はん|たい}の\R{数}{かず}}とよびます。
\end{minipage}

\par\vspace{2mm}
$2-(-3)$ と $2+3$ は、どちらも「2から\R{右}{みぎ}へ3」で5。$2-6$ と $2+(-6)$ は、どちらも「2から\R{左}{ひだり}へ6」で $-4$。だから、\textbf{ひくことは、\R{反対}{はん|たい}の\R{数}{かず}をたすことと\R{同}{おな}じ}です。
\[ 2-(-3)=2+3=5 \qquad 2-6=2+(-6)=-4 \]

\bksubsection{7}{かわった\R{量}{りょう}は「あと－まえ」}

\leavevmode\R{朝}{あさ} \qty{-3}{\degreeCelsius} だった\R{気温}{き|おん}が、\R{昼}{ひる}に \qty{5}{\degreeCelsius} になりました。\R{何}{なん}℃ \R{上}{あ}がったかな？\quad かわった\R{量}{りょう}は「あと－まえ」で\R{求}{もと}めます（1-5）。
\[ 5-(-3)=5+3=8 \]
\noindent\begin{minipage}[c]{0.62\linewidth}
\R{答}{こた}えは $+8$ で、\qty{8}{\degreeCelsius} \R{上}{あ}がりました。\R{温度計}{おん|ど|けい}でたしかめると、$-3$ から0まで3\R{目}{め}もり、0から5まで5\R{目}{め}もりで、\R{合}{あ}わせて8\R{目}{め}もり\R{上}{あ}がっています。

\R{反対}{はん|たい}に、\qty{2}{\degreeCelsius} から \qty{-6}{\degreeCelsius} になったときは、\\ \hspace*{1\zw}$(-6)-2=(-6)+(-2)=-8$\\\R{答}{こた}えのマイナスは「\R{下}{さ}がった」という\R{反対}{はん|たい}の\R{向}{む}きを\R{表}{あらわ}し、\qty{8}{\degreeCelsius} \R{下}{さ}がったことがわかります。
\end{minipage}\hfill
\begin{minipage}[c]{0.35\linewidth}\centering
\begin{tikzpicture}[x=1mm, y=1mm]
  \bkiicthermo{(8,0)}{2.4}{-5}{6}{5}
  \draw[bkOrange, line width=0.6pt, {Stealth[length=3pt]}-{Stealth[length=3pt]}] (16,-7.2) -- (16,0);
  \draw[bkOrange!60!black, line width=0.6pt, {Stealth[length=3pt]}-{Stealth[length=3pt]}] (24,0) -- (24,12);
  \draw[bkGray, line width=0.4pt, densely dashed] (13,-7.2) -- (17,-7.2) (13,0) -- (25,0) (13,12) -- (25,12);
  \node[anchor=west, font=\sffamily\footnotesize, inner sep=0pt] at (17.5,-3.6) {3};
  \node[anchor=west, font=\sffamily\footnotesize, inner sep=0pt] at (25.5,6) {5};
\end{tikzpicture}
\end{minipage}

\bksubsection{8}{かけ\R{算}{ざん}：×$(-1)$ で\R{向}{む}きが\R{反対}{はん|たい}になる}

\noindent{\sffamily\bfseries\color{bkDeep}①\enspace\R{負}{ふ}の\R{数}{かず}×\R{正}{せい}の\R{数}{かず}}\par
かけ\R{算}{ざん}は「1つ\R{分}{ぶん}の\R{数}{かず}×いくつ\R{分}{ぶん}」でした（1-12）。$(-2)\times3$ は「$-2$ が3つ\R{分}{ぶん}」なので、$(-2)+(-2)+(-2)$ です。0から\R{左}{ひだり}へ2ずつ、3\R{回}{かい}\R{進}{すす}んで $-6$ です。

\begin{bkfigure}
\begin{tikzpicture}[x=1mm, y=1mm]
  \bkiicline{0}
  \foreach \s in {0,-2,-4} {\draw[bk arrow, bkBlue, line width=1.1pt] (\s*9,2.5) to[bend right=50] (\s*9-18,2.5);
    \node[anchor=south, font=\footnotesize, text=bkBlue!80!black, inner sep=0pt] at (\s*9-9,9) {$-2$};}
  \bkiicdot{-6}{0}{bkBlue}
\end{tikzpicture}
\bkfigcaption{$(-2)\times3$ は、$-2$ が3つ\R{分}{ぶん}。\R{左}{ひだり}へ2ずつ3\R{回}{かい}で $-6$。}
\end{bkfigure}

\noindent{\sffamily\bfseries\color{bkDeep}②\enspace\R{正}{せい}の\R{数}{かず}×\R{負}{ふ}の\R{数}{かず}}\par
$3\times(-2)$ は「3が $-2$ こ\R{分}{ぶん}」ですが、これでは\R{意味}{い|み}がよくわかりません。そこで、かける\R{数}{かず}を1ずつ\R{小}{ちい}さくしたならびで\R{考}{かんが}えます。

\begin{bkcenter}
\begin{bktable}{colspec={Q[l,m,wd=12mm] *{6}{Q[c,m,wd=14mm]}}, rowsep=1pt, colsep=2.5pt, cells={font=\small}, column{1}={font=\sffamily\bfseries\small}}
\R{式}{しき} & $3\times3$ & $3\times2$ & $3\times1$ & $3\times0$ & $3\times(-1)$ & $3\times(-2)$ \\
\R{答}{こた}え & 9 & 6 & 3 & 0 & {\color{bkRed}\bfseries $-3$} & {\color{bkRed}\bfseries $-6$} \\
\end{bktable}
\bktabcaption{かける\R{数}{かず}が1\R{小}{ちい}さくなると、\R{答}{こた}えは3\R{小}{ちい}さくなる。}
\end{bkcenter}

かける\R{数}{かず}が1へるたびに、3が1つ\R{分}{ぶん}へるので、\R{答}{こた}えは3ずつ\R{小}{ちい}さくなります。このならびをつづけると、$3\times(-1)=-3$、$3\times(-2)=-6$ です。$3\times(-2)$ と $(-2)\times3$ は、どちらも $-6$。かけ\R{算}{ざん}は\R{入}{い}れかえても\R{答}{こた}えが\R{同}{おな}じ（1-13）なので、\R{合}{あ}っています。

$3\times(-1)=-3$ を\R{見}{み}てください。3に $-1$ をかけると、3の\R{反対}{はん|たい}の\R{数}{かず} $-3$ になりました。\textbf{×$(-1)$ は、\R{向}{む}きを\R{反対}{はん|たい}にする}かけ\R{算}{ざん}です。\R{反対}{はん|たい}に\R{言}{い}うと、$-3$ は $3\times(-1)$ のこと。\R{負}{ふ}の\R{数}{かず}は、\R{正}{せい}の\R{数}{かず}に $-1$ をかけた\R{数}{かず}です（$-2=2\times(-1)$、$-5=5\times(-1)$）。

\begin{bkfigure}
\begin{tikzpicture}[x=1mm, y=1mm]
  \bkiicline{0}
  \draw[bk arrow, bkOrange, line width=1.4pt] (0,5) -- (27,5);
  \draw[bk arrow, bkBlue, line width=1.4pt] (0,11) -- (-27,11);
  \draw[bk arrow, bkOrange, line width=1.4pt] (0,17) -- (27,17);
  \node[anchor=west, font=\small, inner sep=0pt] at (30,5) {$3$};
  \node[anchor=east, font=\small, inner sep=0pt] at (-30,11) {$3\times(-1)=-3$};
  \node[anchor=west, font=\small, inner sep=0pt] at (30,17) {$3\times(-1)\times(-1)=3$};
\end{tikzpicture}
\bkfigcaption{×$(-1)$ で\R{向}{む}きが\R{反対}{はん|たい}になる。2\R{回}{かい}で、もとの\R{向}{む}きにもどる。}
\end{bkfigure}

\noindent{\sffamily\bfseries\color{bkDeep}③\enspace\R{負}{ふ}の\R{数}{かず}×\R{負}{ふ}の\R{数}{かず}}\par
$(-2)\times(-3)$ はいくつでしょう。3つの\R{方法}{ほう|ほう}でたしかめます。

\textbf{（ア）\R{時刻}{じ|こく}と\R{場所}{ば|しょ}で}\enspace ミナが\R{西}{にし}へ、1\R{秒}{びょう}に \qty{2}{m} ずつ\R{走}{はし}っています。\R{東}{ひがし}が＋なので、1\R{秒}{びょう}に $-2$\,m ずつ\R{進}{すす}みます。\R{今}{いま}（0\R{秒}{びょう}）は0の\R{場所}{ば|しょ}にいます。3\R{秒後}{びょう|ご}は $(-2)\times3=-6$ で、\R{西}{にし}へ \qty{6}{m} の\R{所}{ところ}です。では、3\R{秒前}{びょう|まえ}はどこにいたかな？\quad「\R{後}{あと}」の\R{反対}{はん|たい}の「\R{前}{まえ}」を\R{負}{ふ}の\R{数}{かず}で\R{表}{あらわ}し、3\R{秒前}{びょう|まえ}を $-3$\R{秒}{びょう}として、\R{表}{ひょう}にします。

\begin{bkcenter}
\begin{bktable}{colspec={Q[l,m,wd=36mm] *{7}{Q[c,m,wd=8.5mm]}}, rowsep=1pt, colsep=4pt, cells={font=\small}, row{1}={font=\sffamily\bfseries\small}}
\R{時刻}{じ|こく}（\R{秒}{びょう}） & $-3$ & $-2$ & $-1$ & 0 & 1 & 2 & 3 \\
\R{場所}{ば|しょ}（m）$(-2)\times$\R{時刻}{じ|こく} & {\color{bkRed}\bfseries 6} & {\color{bkRed}\bfseries 4} & {\color{bkRed}\bfseries 2} & 0 & $-2$ & $-4$ & $-6$ \\
\end{bktable}
\bktabcaption{\R{西}{にし}へ\R{走}{はし}る\R{人}{ひと}は、\R{前}{まえ}の\R{時刻}{じ|こく}には\R{東}{ひがし}（＋）にいた。}
\end{bkcenter}

\R{西}{にし}へ\R{進}{すす}んでいる\R{人}{ひと}は、\R{前}{まえ}の\R{時刻}{じ|こく}には、\R{今}{いま}より\R{東}{ひがし}にいたはずです。1\R{秒前}{びょう|まえ}は\R{東}{ひがし}へ \qty{2}{m}、3\R{秒前}{びょう|まえ}は\R{東}{ひがし}へ \qty{6}{m}。だから3\R{秒前}{びょう|まえ}は $+6$ で、$(-2)\times(-3)=6$ です。

\textbf{（イ）ならびで}\enspace ②と\R{同}{おな}じように、かける\R{数}{かず}を1ずつ\R{小}{ちい}さくします。

\begin{bkcenter}
\begin{bktable}{colspec={Q[l,m,wd=9mm] *{6}{Q[c,m,wd=18.5mm]}}, rowsep=1pt, colsep=2pt, cells={font=\footnotesize}, column{1}={font=\sffamily\bfseries\small}}
\R{式}{しき} & $(-2)\times2$ & $(-2)\times1$ & $(-2)\times0$ & $(-2)\times(-1)$ & $(-2)\times(-2)$ & $(-2)\times(-3)$ \\
\R{答}{こた}え & $-4$ & $-2$ & 0 & {\color{bkRed}\bfseries 2} & {\color{bkRed}\bfseries 4} & {\color{bkRed}\bfseries 6} \\
\end{bktable}
\bktabcaption{かける\R{数}{かず}が1\R{小}{ちい}さくなると、\R{答}{こた}えは2\R{大}{おお}きくなる。}
\end{bkcenter}

$-2$ が1つ\R{分}{ぶん}へるので、\R{答}{こた}えは2ずつ\R{大}{おお}きくなります。このならびをつづけていくと、いちばん\R{右}{みぎ}は $(-2)\times(-3)=6$ になります。

\textbf{（ウ）×$(-1)$ で}\enspace $-2=2\times(-1)$、$-3=3\times(-1)$ なので、
\begin{align*}
(-2)\times(-3)&=2\times(-1)\times3\times(-1) &&\text{（$-2$ と $-3$ を\R{書}{か}きかえる）}\\
&=2\times3\times(-1)\times(-1) &&\text{（かけ\R{算}{ざん}の\R{順番}{じゅん|ばん}を\R{入}{い}れかえる）}\\
&=6\times(-1)\times(-1) &&\text{（6を×$(-1)$ で $-6$、もう1\R{回}{かい}で6）}\\
&=6
\end{align*}
×$(-1)$ を2\R{回}{かい}すると、\R{向}{む}きが2\R{回}{かい}\R{反対}{はん|たい}になって、もとの\R{向}{む}き（＋）にもどります。3つの\R{方法}{ほう|ほう}で、\R{同}{おな}じ6になりました。

\noindent\begin{minipage}[c]{0.46\linewidth}\centering
\begin{bktable}{colspec={Q[c,m,wd=20mm] Q[c,m,wd=12mm] Q[c,m,wd=12mm]}, rowsep=1pt, cells={font=\small}, row{1}={font=\sffamily\bfseries\small}, column{1}={font=\sffamily\bfseries\small}}
かけられる\R{数}{かず} & ×\R{正}{せい} & ×\R{負}{ふ} \\
\R{正}{せい} & ＋ & － \\
\R{負}{ふ} & － & ＋ \\
\end{bktable}
\end{minipage}\hfill
\begin{minipage}[c]{0.51\linewidth}
\R{数}{かず}の\R{前}{まえ}の＋や－を\textbf{\R{符号}{ふ|ごう}}といいます。\bkidx{ふごう}{\R{符号}{ふ|ごう}}かけ\R{算}{ざん}の\R{答}{こた}えの\R{符号}{ふ|ごう}は、\R{同}{おな}じ\R{符号}{ふ|ごう}どうしなら＋、ちがう\R{符号}{ふ|ごう}なら－です。
\end{minipage}

\par\vspace{2mm}
3つ\R{以上}{い|じょう}かけるときは、\R{負}{ふ}の\R{数}{かず}の1つ1つが「×$(-1)$」を1\R{回}{かい}ずつ\R{運}{はこ}んでくると\R{考}{かんが}えます。×$(-1)$ が2\R{回}{かい}・4\R{回}{かい}（\R{偶数}{ぐう|すう}\R{回}{かい}）ならもとの\R{向}{む}きで＋、1\R{回}{かい}・3\R{回}{かい}（\R{奇数}{き|すう}\R{回}{かい}）なら\R{反対}{はん|たい}で－です。

\bksubsection{9}{わり\R{算}{ざん}も、かけ\R{算}{ざん}の\R{逆}{ぎゃく}で}

\leavevmode わり\R{算}{ざん}はかけ\R{算}{ざん}の\R{逆}{ぎゃく}でした（1-14）。$6\div2$ は「2に\R{何}{なに}をかけると6になるか」で、3です。\R{負}{ふ}の\R{数}{かず}でも\R{同}{おな}じように\R{考}{かんが}えます。

\begin{bkcenter}
\begin{bktable}{colspec={Q[c,m,wd=26mm] Q[l,m,wd=62mm] Q[c,m,wd=16mm]}, rowsep=1pt, cells={font=\small}, row{1}={font=\sffamily\bfseries\small}}
わり\R{算}{ざん} & \R{考}{かんが}え\R{方}{かた} & \R{答}{こた}え \\
$6\div2$ & $2\times3=6$ なので & 3 \\
$(-6)\div2$ & $2\times(-3)=-6$ なので & $-3$ \\
$6\div(-2)$ & $(-2)\times(-3)=6$ なので & $-3$ \\
$(-6)\div(-2)$ & $(-2)\times3=-6$ なので & 3 \\
\end{bktable}
\bktabcaption{わる\R{数}{かず}に\R{何}{なに}をかけると、わられる\R{数}{かず}になるか。}
\end{bkcenter}

わり\R{算}{ざん}の\R{答}{こた}えの\R{符号}{ふ|ごう}も、かけ\R{算}{ざん}と\R{同}{おな}じです。\R{同}{おな}じ\R{符号}{ふ|ごう}どうしなら＋、ちがう\R{符号}{ふ|ごう}なら－。0でわることは、\R{負}{ふ}の\R{数}{かず}でもできません（1-14）。

\begin{taisetsu}
マイナスは\R{反対}{はん|たい}の\R{向}{む}き。＋の\R{向}{む}きは、はじめに\R{決}{き}めて\R{変}{か}えない。\\
$a>0$ は\R{正}{せい}、$a<0$ は\R{負}{ふ}。$0<x<6$ は「0より\R{大}{おお}きく6より\R{小}{ちい}さい」。\\
ひくことは、\R{反対}{はん|たい}の\R{数}{かず}をたすことと\R{同}{おな}じ。\\
かけ\R{算}{ざん}・わり\R{算}{ざん}：\R{同}{おな}じ\R{符号}{ふ|ごう}なら＋、ちがう\R{符号}{ふ|ごう}なら－。
\end{taisetsu}

\begin{kotoba}
\small
\bkword{\R{正}{せい}の\R{数}{かず}・\R{負}{ふ}の\R{数}{かず}}{0より\R{大}{おお}きい\R{数}{かず}・0より\R{小}{ちい}さい\R{数}{かず}。0はどちらでもない。}
\bkword{\R{不等号}{ふ|とう|ごう}}{\R{大}{おお}きさをくらべる\R{記号}{き|ごう}。$<$、$>$、$\leqq$、$\geqq$。}
\bkword{\R{整数}{せい|すう}}{…、$-2$、$-1$、0、1、2、…のような\R{数}{かず}。}
\bkword{\R{反対}{はん|たい}の\R{数}{かず}}{0をはさんで\R{反対}{はん|たい}がわの、0から\R{同}{おな}じ\R{距離}{きょ|り}にある\R{数}{かず}。3と $-3$。}
\bkword{\R{符号}{ふ|ごう}}{\R{数}{かず}の\R{前}{まえ}につける＋と－。}
\end{kotoba}

%% ---------------- 例題 ----------------
\bksection{reidai}{\R{例題}{れい|だい}でたしかめよう}

\begin{reidai}[\R{正}{せい}の\R{数}{かず}・\R{負}{ふ}の\R{数}{かず}で\R{表}{あらわ}す]
\R{東}{ひがし}へ \qty{300}{m} \R{進}{すす}むことを $+\qty{300}{m}$ と\R{表}{あらわ}します。（1）\R{西}{にし}へ \qty{200}{m} \R{進}{すす}むことは？\quad（2）$\qty{-500}{m}$ はどんなことを\R{表}{あらわ}しますか。\quad（3）はじめから\R{西}{にし}を＋と\R{決}{き}めていたら、\R{東}{ひがし}へ \qty{300}{m} \R{進}{すす}むことは？\par
\kangaekata ＋と\R{決}{き}めた\R{向}{む}きの\R{反対}{はん|たい}は、マイナスで\R{表}{あらわ}します。\par
\shiki （1）\R{東}{ひがし}が＋なので、\R{反対}{はん|たい}の\R{西}{にし}\R{向}{む}きはマイナス。\par
\hspace*{2.6\zw}（2）マイナスなので、＋（\R{東}{ひがし}）の\R{反対}{はん|たい}の\R{西}{にし}\R{向}{む}き。\par
\hspace*{2.6\zw}（3）\R{西}{にし}が＋なので、\R{反対}{はん|たい}の\R{東}{ひがし}\R{向}{む}きはマイナス。\par
\kotae （1）$\qty{-200}{m}$\quad（2）\R{西}{にし}へ \qty{500}{m} \R{進}{すす}むこと\quad（3）$\qty{-300}{m}$
\end{reidai}

\begin{jibun}
5\,kg ふえることを $+\qty{5}{kg}$ と\R{表}{あらわ}します。（1）2\,kg へることは？\quad（2）$\qty{-3}{kg}$ はどんなことを\R{表}{あらわ}しますか。\quad（3）はじめから「へる」を＋と\R{決}{き}めていたら、5\,kg ふえることは？
\bkjans{（1）\nobreak $\qty{-2}{kg}$\quad（2）\nobreak 3\,kg へること\quad（3）\nobreak $\qty{-5}{kg}$}
\end{jibun}

\begin{reidai}[\R{不等号}{ふ|とう|ごう}と\R{以上}{い|じょう}・\R{未満}{み|まん}]
（1）$-5$ と $-2$、（2）$-1.5$ と $0.5$ の\R{大}{おお}きさを、\R{不等号}{ふ|とう|ごう}で\R{表}{あらわ}しましょう。（3）$-3$ \R{以上}{い|じょう} 2 \R{未満}{み|まん}の\R{整数}{せい|すう}を、すべて\R{書}{か}きましょう。\par
\kangaekata \R{数直線}{すう|ちょく|せん}で、\R{左}{ひだり}にあるほうが\R{小}{ちい}さい。\R{以上}{い|じょう}はその\R{数}{かず}をふくみ、\R{未満}{み|まん}はふくまない。\par
\begin{bkcenter}
\begin{tikzpicture}[x=1mm, y=1mm]
  \bkiicline{0}
  \bkiicdot{-5}{0}{bkBlue}\bkiicdot{-2}{0}{bkBlue}\bkiicdot{-1.5}{0}{bkOrange}\bkiicdot{0.5}{0}{bkOrange}
  \foreach \v/\l in {-5/{$-5$}, -2/{$-2$}} {\node[anchor=south, font=\footnotesize, text=bkBlue!80!black, inner sep=0pt] at (\v*9,2.6) {\l};}
  \foreach \v/\l in {-1.5/{$-1.5$}, 0.5/{$0.5$}} {\node[anchor=south, font=\footnotesize, text=bkOrange!80!black, inner sep=0pt] at (\v*9,2.6) {\l};}
\end{tikzpicture}
\end{bkcenter}
\shiki （1）$-5$ は $-2$ より\R{左}{ひだり}にある。\quad（2）$-1.5$ は $0.5$ より\R{左}{ひだり}にある。\par
\hspace*{2.6\zw}（3）$-3$ はふくみ、2 はふくまない。$-3$ から1までの\R{整数}{せい|すう}。\par
\kotae （1）$-5<-2$\quad（2）$-1.5<0.5$\quad（3）$-3$、$-2$、$-1$、0、1
\end{reidai}

\begin{jibun}
（1）$-1$ と $-4$、（2）$0$ と $-2.5$ の\R{大}{おお}きさを、\R{不等号}{ふ|とう|ごう}で\R{表}{あらわ}しましょう。（3）$-2$ より\R{大}{おお}きく3\R{以下}{い|か}の\R{整数}{せい|すう}を、すべて\R{書}{か}きましょう。
\bkjans{（1）\nobreak $-1>-4$（$-4<-1$）\quad（2）\nobreak $0>-2.5$（$-2.5<0$）\quad（3）\nobreak $-1$、0、1、2、3}
\end{jibun}

\begin{reidai}[たし\R{算}{ざん}]
（1）$(-3)+5$\quad（2）$2+(-6)$\quad（3）$(-3)+(-2)$\par
\kangaekata \R{正}{せい}の\R{数}{かず}をたすと\R{右}{みぎ}へ、\R{負}{ふ}の\R{数}{かず}をたすと\R{左}{ひだり}へ\R{進}{すす}みます。\par
\shiki （1）$-3$ から\R{右}{みぎ}へ5\R{進}{すす}む。$-3$ から0まで3、さらに2で、2。\par
\hspace*{2.6\zw}（2）2から\R{左}{ひだり}へ6\R{進}{すす}む。2から0まで2、さらに4で、$-4$。\par
\hspace*{2.6\zw}（3）$-3$ から\R{左}{ひだり}へ2\R{進}{すす}んで、$-5$。\par
\kotae （1）2\quad（2）$-4$\quad（3）$-5$
\end{reidai}

\begin{jibun}
（1）$(-4)+7$\quad（2）$3+(-5)$\quad（3）$(-1)+(-6)$
\bkjans{（1）\nobreak 3\quad（2）\nobreak $-2$\quad（3）\nobreak $-7$}
\end{jibun}

\begin{reidai}[ひき\R{算}{ざん}]
（1）$(-1)-4$\quad（2）$3-(-2)$\quad（3）$(-2)-(-6)$\par
\kangaekata ひき\R{算}{ざん}は、\R{反対}{はん|たい}の\R{数}{かず}をたすたし\R{算}{ざん}になおしてから、\R{数直線}{すう|ちょく|せん}で\R{進}{すす}みます。\par
\begin{bkcenter}
\begin{tikzpicture}[x=1mm, y=1mm]
  \bkiicline{0}
  \bkiicmove{-1}{-5}{5}{bkBlue}{（1）}
  \bkiicmove{3}{5}{5}{bkOrange}{（2）}
  \bkiicmove{-2}{4}{13}{bkOrange}{（3）}
\end{tikzpicture}
\end{bkcenter}
\shiki （1）$(-1)-4=(-1)+(-4)$。$-1$ から\R{左}{ひだり}へ4で、$-5$。\par
\hspace*{2.6\zw}（2）$3-(-2)=3+2$。3から\R{右}{みぎ}へ2で、5。\par
\hspace*{2.6\zw}（3）$(-2)-(-6)=(-2)+6$。$-2$ から\R{右}{みぎ}へ6で、4。\par
\kotae （1）$-5$\quad（2）5\quad（3）4
\end{reidai}

\begin{jibun}
（1）$1-5$\quad（2）$(-4)-(-7)$\quad（3）$0-(-3)$
\bkjans{（1）\nobreak $-4$\quad（2）\nobreak 3\quad（3）\nobreak 3}
\end{jibun}

\begin{reidai}[かわった\R{量}{りょう}と\R{差}{さ}]
（1）\qty{-3}{\degreeCelsius} → \qty{5}{\degreeCelsius} の\R{変化}{へん|か}\quad（2）\qty{2}{\degreeCelsius} → \qty{-6}{\degreeCelsius} の\R{変化}{へん|か}\par
（3）\R{最高気温}{さい|こう|き|おん}が \qty{-4}{\degreeCelsius}、\R{最低気温}{さい|てい|き|おん}が \qty{-9}{\degreeCelsius} の\R{日}{ひ}の、\R{最高}{さい|こう}と\R{最低}{さい|てい}の\R{差}{さ}\par
\kangaekata かわった\R{量}{りょう}は「あと－まえ」。\R{差}{さ}は「\R{大}{おお}きいほう－\R{小}{ちい}さいほう」です。\par
\begin{bkcenter}
\begin{tikzpicture}[x=1mm, y=1mm]
  % （3）の数直線（1 目もり 10 mm）
  \draw[bkInk, line width=0.9pt, -{Stealth[length=5pt]}] (-103,0) -- (12,0);
  \foreach \k in {-10,...,1} {\draw[bkInk, line width=0.6pt] (\k*10,-1.3) -- (\k*10,1.3);
    \node[anchor=north, font=\footnotesize, inner sep=0pt] at (\k*10,-2.6) {$\k$};}
  \fill[bkBlue] (-90,0) circle (1.5);\fill[bkOrange] (-40,0) circle (1.5);
  \node[anchor=south, font=\sffamily\footnotesize, text=bkBlue!80!black, inner sep=0pt] at (-90,3) {\R{最低}{さい|てい}};
  \node[anchor=south, font=\sffamily\footnotesize, text=bkOrange!80!black, inner sep=0pt] at (-40,3) {\R{最高}{さい|こう}};
  \draw[bkPurple, line width=0.7pt, {Stealth[length=3pt]}-{Stealth[length=3pt]}] (-88.5,12) -- (-41.5,12);
  \node[anchor=south, font=\sffamily\footnotesize, text=bkPurple, inner sep=0pt] at (-65,13.5) {\R{差}{さ}は5\R{目}{め}もり};
  \node[anchor=west, font=\sffamily\footnotesize, text=bkGray, inner sep=0pt] at (13,0) {℃};
\end{tikzpicture}
\end{bkcenter}
\shiki （1）$5-(-3)=5+3=8$\par
\hspace*{2.6\zw}（2）$(-6)-2=(-6)+(-2)=-8$\par
\hspace*{2.6\zw}（3）$-4$ のほうが\R{大}{おお}きい。$(-4)-(-9)=(-4)+9=5$\par
\kotae （1）$+8$\,℃（\qty{8}{\degreeCelsius} \R{上}{あ}がった）\quad（2）$-8$\,℃（\qty{8}{\degreeCelsius} \R{下}{さ}がった）\quad（3）\qty{5}{\degreeCelsius}
\end{reidai}

\begin{jibun}
（1）\qty{-7}{\degreeCelsius} → \qty{-1}{\degreeCelsius} の\R{変化}{へん|か}\quad（2）\qty{4}{\degreeCelsius} → \qty{-3}{\degreeCelsius} の\R{変化}{へん|か}\par
（3）\R{最高気温}{さい|こう|き|おん}が \qty{1}{\degreeCelsius}、\R{最低気温}{さい|てい|き|おん}が \qty{-6}{\degreeCelsius} の\R{日}{ひ}の、\R{最高}{さい|こう}と\R{最低}{さい|てい}の\R{差}{さ}
\bkjans{（1）\nobreak $(-1)-(-7)=6$。\qty{6}{\degreeCelsius} \R{上}{あ}がった\quad（2）\nobreak $(-3)-4=-7$。\qty{7}{\degreeCelsius} \R{下}{さ}がった\quad（3）\nobreak $1-(-6)=7$。\qty{7}{\degreeCelsius}}
\end{jibun}

\begin{reidai}[かけ\R{算}{ざん}]
（1）$(-4)\times3$\quad（2）$5\times(-2)$\quad（3）$(-3)\times(-4)$\quad（4）$(-2)\times(-3)\times(-1)$\par
\kangaekata \R{符号}{ふ|ごう}をのぞいて\R{計算}{けい|さん}します。\R{負}{ふ}の\R{数}{かず}1つが×$(-1)$ 1\R{回}{かい}なので、その\R{回数}{かい|すう}で\R{符号}{ふ|ごう}を\R{決}{き}めます。\par
\shiki （1）$4\times3=12$。\R{負}{ふ}の\R{数}{かず}が1つ（×$(-1)$ が1\R{回}{かい}）で－\par
\hspace*{2.6\zw}（2）$5\times2=10$。×$(-1)$ が1\R{回}{かい}で－\par
\hspace*{2.6\zw}（3）$3\times4=12$。×$(-1)$ が2\R{回}{かい}で、もとの\R{向}{む}き（＋）\par
\hspace*{2.6\zw}（4）$2\times3\times1=6$。×$(-1)$ が3\R{回}{かい}で、\R{反対}{はん|たい}の\R{向}{む}き（－）\par
\kotae （1）$-12$\quad（2）$-10$\quad（3）12\quad（4）$-6$
\end{reidai}

\begin{jibun}
（1）$(-6)\times2$\quad（2）$3\times(-7)$\quad（3）$(-5)\times(-2)$\quad（4）$(-1)\times4\times(-3)$
\bkjans{（1）\nobreak $-12$\quad（2）\nobreak $-21$\quad（3）\nobreak 10\quad（4）\nobreak 12}
\end{jibun}

\begin{reidai}[わり\R{算}{ざん}]
（1）$(-12)\div(-3)$\quad（2）$(-15)\div3$\quad（3）$16\div(-4)$\par
\kangaekata 「わる\R{数}{かず}に\R{何}{なに}をかけると、わられる\R{数}{かず}になるか」を\R{考}{かんが}えます。\R{符号}{ふ|ごう}は、かけ\R{算}{ざん}と\R{同}{おな}じきまりです。\par
\shiki （1）$(-3)\times4=-12$ なので4。\R{同}{おな}じ\R{符号}{ふ|ごう}どうしで＋\par
\hspace*{2.6\zw}（2）$3\times(-5)=-15$ なので $-5$。ちがう\R{符号}{ふ|ごう}で－\par
\hspace*{2.6\zw}（3）$(-4)\times(-4)=16$ なので $-4$。ちがう\R{符号}{ふ|ごう}で－\par
\kotae （1）4\quad（2）$-5$\quad（3）$-4$
\end{reidai}

\begin{jibun}
（1）$(-18)\div6$\quad（2）$(-24)\div(-8)$\quad（3）$21\div(-7)$
\bkjans{（1）\nobreak $-3$\quad（2）\nobreak 3\quad（3）\nobreak $-3$}
\end{jibun}

\begin{reidai}[\R{向}{む}きのある\R{量}{りょう}のかけ\R{算}{ざん}]
\R{東}{ひがし}を＋とします。\R{西}{にし}へ1\R{秒}{びょう}に \qty{3}{m} ずつゆっくり\R{進}{すす}む\R{車}{くるま}が、\R{今}{いま}0の\R{場所}{ば|しょ}にいます。2\R{秒後}{びょう|ご}と2\R{秒前}{びょう|まえ}の\R{場所}{ば|しょ}は？\par
\kangaekata 1\R{秒}{びょう}に $-3$\,m。\R{時刻}{じ|こく}は、2\R{秒後}{びょう|ご}が $+2$、2\R{秒前}{びょう|まえ}が $-2$ です。\R{場所}{ば|しょ}＝（1\R{秒}{びょう}に\R{進}{すす}む\R{量}{りょう}）×（\R{時刻}{じ|こく}）。\par
\shiki 2\R{秒後}{びょう|ご}：$(-3)\times2=-6$\qquad 2\R{秒前}{びょう|まえ}：$(-3)\times(-2)=6$\par
\begin{bkcenter}
\begin{tikzpicture}[x=1mm, y=1mm]
  \node[anchor=south west, font=\sffamily\small, inner sep=0pt] at (-58.5,2.5) {\R{西}{にし}};
  \node[anchor=south east, font=\sffamily\small, inner sep=0pt] at (60.5,2.5) {\R{東}{ひがし}（＋）};
  \bkiicline{0}
  \bkiicdot{-6}{0}{bkBlue}\bkiicdot{0}{0}{bkInk}\bkiicdot{6}{0}{bkOrange}
  \node[anchor=south, font=\sffamily\footnotesize, text=bkBlue!80!black, inner sep=0pt] at (-54,9) {2\R{秒後}{びょう|ご}};
  \node[anchor=south, font=\sffamily\footnotesize, inner sep=0pt] at (0,9) {\R{今}{いま}};
  \node[anchor=south, font=\sffamily\footnotesize, text=bkOrange!80!black, inner sep=0pt] at (54,9) {2\R{秒前}{びょう|まえ}};
  \draw[bk arrow, bkGray, line width=1.1pt] (36,5) -- node[above, font=\sffamily\footnotesize, text=bkInk, inner sep=1mm] {\R{進}{すす}む\R{向}{む}き} (18,5);
\end{tikzpicture}
\end{bkcenter}
\kotae 2\R{秒後}{びょう|ご}は\R{西}{にし}へ \qty{6}{m}、2\R{秒前}{びょう|まえ}は\R{東}{ひがし}へ \qty{6}{m}
\end{reidai}

\begin{jibun}
\R{上}{うえ}が＋。1\R{秒}{びょう}に \qty{2}{m} ずつ\R{下}{さ}がるエレベーターが、\R{今}{いま}\R{高}{たか}さ0の\R{所}{ところ}にいます。3\R{秒後}{びょう|ご}と3\R{秒前}{びょう|まえ}の\R{高}{たか}さは？
\bkjans{\mbox{3\R{秒後}{びょう|ご}}：$(-2)\times3=-6$ で $-6$\,m\quad \mbox{3\R{秒前}{びょう|まえ}}：$(-2)\times(-3)=6$ で $+6$\,m}
\end{jibun}

%% ---------------- まちがえやすいところ・物理工学 ----------------
\bksection{tsumazuki}{まちがえやすいところをチェック}

\begin{tsumazuki}{\R{負}{ふ}の\R{数}{かず}の\R{大}{おお}きさ}
\bkng{$-5>-2$\quad（5は2より\R{大}{おお}きいから）}
\bkok{$-5<-2$\quad（$-5$ のほうが\R{左}{ひだり}。\qty{-5}{\degreeCelsius} のほうが\R{寒}{さむ}い）}
\R{負}{ふ}の\R{数}{かず}は、0から\R{遠}{とお}いほど\R{小}{ちい}さくなります。
\end{tsumazuki}

\begin{tsumazuki}{\R{負}{ふ}の\R{数}{かず}をひく}
\bkng{$4-(-3)=1$}
\bkok{$4-(-3)=4+3=7$\quad（$-3$ をひく＝\R{反対}{はん|たい}の $+3$ をたす）}
「$-$」が2つ\R{続}{つづ}いたら、ひき\R{算}{ざん}を\R{反対}{はん|たい}の\R{数}{かず}のたし\R{算}{ざん}に\R{書}{か}きなおしてから\R{計算}{けい|さん}します。
\end{tsumazuki}

\begin{tsumazuki}{\R{負}{ふ}の\R{数}{かず}から\R{正}{せい}の\R{数}{かず}をひく}
\bkng{$(-3)-2=-1$\quad（3から2をひいて1にした）}
\bkok{$(-3)-2=(-3)+(-2)=-5$\quad（$-3$ から\R{左}{ひだり}へ2）}
ひくと\R{左}{ひだり}へ\R{進}{すす}むので、\R{答}{こた}えは $-3$ より\R{小}{ちい}さくなるはずです。
\end{tsumazuki}

\begin{tsumazuki}{\R{負}{ふ}の\R{数}{かず}どうしのかけ\R{算}{ざん}}
\bkng{$(-3)\times(-4)=-12$\quad（マイナスがあるから－）}
\bkok{$(-3)\times(-4)=12$\quad（×$(-1)$ が2\R{回}{かい}で、もとの\R{向}{む}き）}
たし\R{算}{ざん}の $(-3)+(-4)=-7$ とまぜないようにしよう。
\end{tsumazuki}

\begin{tsumazuki}{とちゅうで＋の\R{向}{む}きを\R{変}{か}える}
\leavevmode\label{bk:2-2:tochu}\bkng{\R{東}{ひがし}へ \qty{3}{m}（$+3$）、\R{次}{つぎ}に\R{西}{にし}へ \qty{5}{m} を $+5$ と\R{書}{か}いて、$3+5=8$}
\bkok{\R{西}{にし}へ \qty{5}{m} は $-5$。$3+(-5)=-2$ で、\R{出発点}{しゅっ|ぱつ|てん}の\R{西}{にし}へ \qty{2}{m}}
\R{東}{ひがし}を＋と\R{決}{き}めたら、\R{最後}{さい|ご}まで\R{東}{ひがし}が＋です。
\end{tsumazuki}

\begin{tsumazuki}{「\R{未満}{み|まん}」に、その\R{数}{かず}を\R{入}{い}れてしまう}
\bkng{6\R{未満}{み|まん}の\R{整数}{せい|すう}に、6を\R{入}{い}れる}
\bkok{6は\R{入}{はい}らない（6\R{以下}{い|か}なら\R{入}{はい}る。$\leqq$ の\R{下}{した}の\R{線}{せん}が「\R{等}{ひと}しい」）}
\end{tsumazuki}

\begin{butsuri}
\noindent\begin{minipage}[c]{0.38\linewidth}\centering
\begin{tikzpicture}[x=1mm, y=1mm]
  % ℃ と K の目もり（1目もり 100 度を 12 mm）
  \foreach \x/\u in {4/℃, 30/K} {\draw[bkInk, line width=0.9pt] (\x,-33.6) -- (\x,13);
    \node[font=\sffamily\small, text=bkInk] at (\x,17) {\u};}
  \foreach \c/\k/\col in {100/373/bkInk, 0/273/bkBlue, -273/0/bkPurple} {
    \pgfmathsetmacro\yy{\c*0.12}
    \draw[bkInk, line width=0.6pt] (2.5,\yy) -- (5.5,\yy) (28.5,\yy) -- (31.5,\yy);
    \draw[bkGray, line width=0.4pt, densely dashed] (7,\yy) -- (27,\yy);
    \node[anchor=east, font=\footnotesize, text=\col, inner sep=0pt] at (1.5,\yy) {$\c$};
    \node[anchor=west, font=\footnotesize, text=\col, inner sep=0pt] at (33,\yy) {$\k$};}
  \fill[bkPurple] (4,-32.76) circle (1.2) (30,-32.76) circle (1.2);
  \node[anchor=north, font=\sffamily\footnotesize, text=bkPurple!85!black, inner sep=0pt] at (17,-36) {いちばん\R{低}{ひく}い\R{温度}{おん|ど}};
\end{tikzpicture}
\end{minipage}\hfill
\begin{minipage}[c]{0.59\linewidth}
\textbf{0\,℃ より\R{下}{した}にも、\R{底}{そこ}がある}\par
\R{約}{やく} \qty{-273}{\degreeCelsius} より\R{低}{ひく}い\R{温度}{おん|ど}はありません。そこで\R{科学}{か|がく}では、このいちばん\R{低}{ひく}い\R{温度}{おん|ど}を0とする\R{目}{め}もり（\R{単位}{たん|い}は K、ケルビン）も\R{使}{つか}います。\R{超電導}{ちょう|でん|どう}リニアには、\R{約}{やく} \qty{-269}{\degreeCelsius}（\R{約}{やく} \qty{4}{K}）まで\R{冷}{ひ}やして\R{使}{つか}う\R{磁石}{じ|しゃく}があります。K の\R{数}{かず}は、℃ の\R{数}{かず}に\R{約}{やく}273をたした\R{数}{かず}です（$-269+273=4$、$0+273=273$）。
\end{minipage}

\par\vspace{2mm}
\textbf{℃の\R{向}{む}き}\par
℃の\R{目}{め}もりを\R{決}{き}めたセルシウスは、はじめ、ふっとうする\R{温度}{おん|ど}を0、こおる\R{温度}{おん|ど}を100にしていました。\R{向}{む}きは\R{決}{き}め\R{方}{かた}しだいなのです。

\par\vspace{2mm}
\noindent\begin{minipage}[c]{0.58\linewidth}
\textbf{\R{飛行機}{ひ|こう|き}が\R{飛}{と}ぶ\R{空}{そら}は、とても\R{寒}{さむ}い}\par
\R{空気}{くう|き}は、\R{高}{たか}い\R{所}{ところ}ほど\R{冷}{つめ}たくなります。\R{旅客機}{りょ|かく|き}が\R{飛}{と}ぶ\R{高}{たか}さ\R{約}{やく} \qty{10000}{m} では、\R{外}{そと}の\R{気温}{き|おん}は\R{約}{やく} \qty{-50}{\degreeCelsius} です。\R{地上}{ち|じょう}が \qty{15}{\degreeCelsius} の\R{日}{ひ}なら、\R{差}{さ}は $15-(-50)=65$ で、\R{約}{やく} \qty{65}{\degreeCelsius}。\R{飛行機}{ひ|こう|き}は、\R{温度計}{おん|ど|けい}（センサー）でこのマイナスの\R{温度}{おん|ど}をはかり、\R{窓}{まど}やつばさが\R{凍}{こお}らないように\R{温}{あたた}めたり、\R{客室}{きゃく|しつ}を\R{暖}{あたた}かく\R{保}{たも}ったりしています。
\end{minipage}\hfill
\begin{minipage}[c]{0.40\linewidth}\centering
\begin{tikzpicture}[x=1mm, y=1mm]
  \draw[bkInk, line width=0.8pt, -{Stealth[length=4pt]}] (0,0) -- (0,40);
  \node[anchor=south, font=\sffamily\footnotesize, inner sep=0pt] at (0,41.5) {\R{高}{たか}さ};
  \foreach \yy/\h in {0/0, 34/10000} {\draw[bkInk, line width=0.6pt] (-1.2,\yy) -- (1.2,\yy);
    \node[anchor=east, font=\footnotesize, inner sep=0pt] at (-2.5,\yy) {\h\,m};}
  \fill[bkStem] (-1,-1.2) rectangle (30,0);
  \begin{scope}[shift={(11,34)}]
    \fill[bkGray!60] (-7,-1.2) rectangle (9,1.2);
    \fill[bkGray!60] (-7,-1.2) -- (-9,3.5) -- (-6,3.5) -- (-4,1.2) -- cycle;
    \fill[bkGray!80] (-1,0) -- (3,-4) -- (5,-4) -- (3,0) -- cycle;
  \end{scope}
  \node[anchor=west, font=\sffamily\footnotesize, text=bkBlue!80!black, inner sep=0pt] at (21,34) {\qty{-50}{\degreeCelsius}};
  \node[anchor=west, font=\sffamily\footnotesize, text=bkOrange!80!black, inner sep=0pt] at (21,4) {\qty{15}{\degreeCelsius}};
\end{tikzpicture}
\end{minipage}

\end{butsuri}



\clearpage
%% ---------------- 章末問題 ----------------
\bkendproblems
\bkproblemgroup{\R{基本}{き|ほん}}{1}

\bkproblem \R{海面}{かい|めん}より \qty{40}{m} \R{高}{たか}い\R{所}{ところ}を $+\qty{40}{m}$ と\R{表}{あらわ}します。（1）\R{海面}{かい|めん}より \qty{15}{m} \R{低}{ひく}い\R{所}{ところ}は？\quad（2）$\qty{-120}{m}$ はどんな\R{所}{ところ}ですか。
\bkpans{（1）\nobreak $\qty{-15}{m}$\quad（2）\nobreak \R{海面}{かい|めん}より \qty{120}{m} \R{低}{ひく}い\R{所}{ところ}}

\bkproblem □に $<$ か $>$ を\R{入}{い}れましょう。\R{下}{した}の\R{数直線}{すう|ちょく|せん}を\R{使}{つか}ってもかまいません。\\
\mbox{（1）$-6$ □ $-1$}\quad \mbox{（2）$0$ □ $-3$}\quad \mbox{（3）$-2.5$ □ $2$}
\begin{bkcenter}
\begin{tikzpicture}[x=1mm, y=1mm]
  \bkiicline{0}
\end{tikzpicture}
\end{bkcenter}
\bkpans{（1）\nobreak $<$\quad（2）\nobreak $>$\quad（3）\nobreak $<$}

\bkproblem （1）$-4$ \R{以上}{い|じょう} 1 \R{以下}{い|か}の\R{整数}{せい|すう}を、すべて\R{書}{か}きましょう。（2）$-2$ より\R{大}{おお}きく 3 \R{未満}{み|まん}の\R{整数}{せい|すう}を、すべて\R{書}{か}きましょう。ノートに\R{下}{した}のような\R{数直線}{すう|ちょく|せん}をかいて、\R{入}{はい}る\R{数}{かず}に\R{印}{しるし}をつけよう。
\begin{bkcenter}
\begin{tikzpicture}[x=1mm, y=1mm]
  \bkiicline{0}
\end{tikzpicture}
\end{bkcenter}
\bkpans{（1）\nobreak $-4$、$-3$、$-2$、$-1$、0、1\quad（2）\nobreak $-1$、0、1、2}

\bkproblem \R{次}{つぎ}の\R{計算}{けい|さん}をしましょう。\\
（1）$(-2)+6$\quad（2）$1+(-4)$\quad（3）$(-3)+(-2)$\\（4）$3-7$\quad（5）$(-5)-(-1)$\quad（6）$0-(-4)$
\bkpans{（1）\nobreak 4\quad（2）\nobreak $-3$\quad（3）\nobreak $-5$\quad（4）\nobreak $-4$\quad（5）\nobreak $-4$\quad（6）\nobreak 4}

\bkproblem \R{次}{つぎ}の\R{計算}{けい|さん}をしましょう。\\
（1）$(-7)\times3$\quad（2）$4\times(-5)$\quad（3）$(-6)\times(-6)$\\（4）$(-20)\div4$\quad（5）$(-15)\div(-5)$
\bkpans{（1）\nobreak $-21$\quad（2）\nobreak $-20$\quad（3）\nobreak 36\quad（4）\nobreak $-5$\quad（5）\nobreak 3}

\bkproblemgroup{チャレンジ}{2}

\noindent\begin{minipage}[c]{0.70\linewidth}
\bkproblem ある\R{日}{ひ}の\R{南極}{なん|きょく}の\R{気温}{き|おん}は、\R{最高}{さい|こう}が \qty{-2}{\degreeCelsius}、\R{最低}{さい|てい}が \qty{-11}{\degreeCelsius} でした。\R{最高}{さい|こう}と\R{最低}{さい|てい}の\R{差}{さ}は\R{何}{なん}℃ ですか。\R{右}{みぎ}の\R{温度計}{おん|ど|けい}で\R{確}{たし}かめよう。
\end{minipage}\hfill
\begin{minipage}[c]{0.27\linewidth}\centering
\begin{tikzpicture}[x=1mm, y=1mm]
  \bkiicthermo{(8,0)}{1.5}{-15}{5}{-2}
  \bkiicthermo{(28,0)}{1.5}{-15}{5}{-11}
  \node[anchor=north, font=\sffamily\footnotesize, inner sep=0pt] at (8,-33) {\R{最高}{さい|こう}};
  \node[anchor=north, font=\sffamily\footnotesize, inner sep=0pt] at (28,-33) {\R{最低}{さい|てい}};
\end{tikzpicture}
\end{minipage}
\bkpans{$(-2)-(-11)=9$。\qty{9}{\degreeCelsius}}

\noindent\begin{minipage}[c]{0.70\linewidth}
\bkproblem \R{上}{うえ}\R{向}{む}きを＋と\R{決}{き}めます。ボールが\R{下}{した}へ \qty{15}{m} \R{落}{お}ちたとき、\R{高}{たか}さのかわった\R{量}{りょう}は\R{何}{なん}m ですか。\R{下}{した}\R{向}{む}きを＋と\R{決}{き}めたら、どうなりますか。
\end{minipage}\hfill
\begin{minipage}[c]{0.27\linewidth}\centering
\begin{tikzpicture}[x=1mm, y=1mm]
  \draw[bkInk, line width=0.9pt, -{Stealth[length=5pt]}] (0,-2) -- (0,30);
  \node[anchor=south, font=\sffamily\small, inner sep=0pt] at (0,31.5) {\R{上}{うえ}（＋）};
  \draw[bkGray, line width=0.6pt, dash pattern=on 2pt off 1.5pt] (10,24) circle (2.6);
  \shade[ball color=bkOrange] (10,2) circle (2.6);
  \draw[bk arrow, bkBlue, line width=1.2pt] (10,20.5) -- (10,6);
  \node[anchor=west, font=\sffamily\small, inner sep=0pt] at (12.5,13) {\qty{15}{m}};
\end{tikzpicture}
\end{minipage}
\bkpans{$\qty{-15}{m}$。\R{下}{した}\R{向}{む}きを＋にすると $+\qty{15}{m}$}

\bkproblem \R{次}{つぎ}の\R{計算}{けい|さん}をしましょう。\quad（1）$(-2)\times3\times(-5)$\quad（2）$(-1)\times(-1)\times(-1)$
\bkpans{（1）\nobreak 30\quad（2）\nobreak $-1$}

\begin{furikaeri}
\begin{checklist}
\item 0より\R{小}{ちい}さい\R{数}{かず}を\R{負}{ふ}の\R{数}{かず}で\R{表}{あらわ}し、\R{数直線}{すう|ちょく|せん}にかける。\hfill{\footnotesize\sffamily\color{bkTheme}→\pageref{bk:2-2:fu}ページ}
\item \R{不等号}{ふ|とう|ごう}（$<$・$>$・$\leqq$・$\geqq$）で、\R{数}{かず}の\R{大}{おお}きさをくらべられる。\hfill{\footnotesize\sffamily\color{bkTheme}→\pageref{bk:2-2:futougou}ページ}
\item \R{正}{せい}の\R{数}{かず}・\R{負}{ふ}の\R{数}{かず}の ＋ － × ÷ の\R{計算}{けい|さん}ができる。\hfill{\footnotesize\sffamily\color{bkTheme}→\pageref{bk:2-2:keisan}ページ}
\end{checklist}
チェックがつかなかったところは、→のページにもどって、もう\R{一度}{いち|ど}\R{読}{よ}んでみよう。
\end{furikaeri}

\bkfukidashi{sora}{マイナスは「\R{反対}{はん|たい}の\R{向}{む}き」。\R{向}{む}きのある\R{量}{りょう}は、この\R{先}{さき}、\R{矢印}{や|じるし}（ベクトル）に\R{育}{そだ}つんだって。}
